%==============================================================================
% MATH 5080 — Introduction to Analysis I (Fall 2026)
% Class:        5
% Date:         Sept 8th
%==============================================================================

\begin{theorem}[Bolzano-Weiestrass]
    If $S \subset \mathbb{R}$ is bounded and contains infinitely many points then $S^{\prime} \neq \emptyset$.
\end{theorem}

\begin{remark}
    Bolzano-Weierstrass gives the existence of limit points, but it does not tell us 
    that the limit points are in $S$. 
\end{remark}

\begin{proof}
    Suppose $S^{\prime} = \emptyset$, then $S = \bar{S} $ is closed. 
    By HB, $S$ is compact.
    \medskip 

    For any $x \in S, x \notin S^{\prime}$, $x$ must be isolated point, i.e. 
    \[
        \exists \ep_x > 0, s.t. N(x, \ep_x) \cap S = \{ x \} 
    \]
    Then $ \{ N(x, \varepsilon_x) \}_{x \in S} $ is an open cover of $S$. 
    By compactness, $\exists x_1, x_2, \hdots x_N$ s.t. 
    \[
        N(x_1, \varepsilon_{x_1}) \cup \hdots N(x_N, \varepsilon_{x_N}) \supset S
    \]
    However, each one of the neighborhoods contains only its center, i.e. 
    \[
        N(x_1, \varepsilon_{x_1}) = \{ x_1 \}, \hdots N(x_n, \varepsilon_{x_N}) = \{ x_N \} 
    \]
    Hence 
    \[
        \{ x_1 \}  \cup \hdots \{ x_N \}  \supset S
    \]
    and this contradicts $S$ having infinitely many points.
\end{proof}

\begin{remark}
    $S$ has infinite points and $S^{\prime} \neq \emptyset \not \implies S$ bounded. 
\end{remark}

\begin{proof}
    Let $S \coloneq [1, 2] \cup \mathbb{N}$.
    Then $S^{\prime} = [1, 2] \neq \emptyset$  , but $S$ unbounded.
\end{proof}

\begin{remark}
    If $S \subset \mathbb{R}$, and every 
    infinite subset of $\mathbb{S}$ has an accumulation point in $\mathbb{R}$, 
    then $S$ is bounded.
\end{remark}

\begin{proof}
    \remind{To be added in}
\end{proof}

\begin{remark}[Relationship between results]
    \hfill
    \begin{center}
    \begin{tikzpicture}[
        thm/.style={
            draw, 
            rounded corners=2pt, 
            fill=blue!10, 
            minimum height=1.5cm,
            minimum width=3.2cm,
            align=center,
            inner sep=4pt
        },
        imp/.style={
            double,
            double distance=1.6pt,
            -{Implies[length=2.6mm]}
        },
        eqv/.style={
            double,
            double distance=1.6pt,
            {Implies[length=2.6mm]}-{Implies[length=2.6mm]}
        },
        lbl/.style={font=\small, fill=white, inner sep=1.5pt}
    ]
        \node[thm] (CA)  at (0,0)    {Completeness\\Axiom (CA)};
        \node[thm] (MCT) at (4.2,0)  {Monotone\\Convergence (MCT)};
        \node[thm] (BW)  at (8.4,0)  {Bolzano--\\Weierstrass (BW)};
        \node[thm] (HB)  at (12.6,0) {Heine--Borel\\(HB)};
        \node[thm] (NIP) at (0,-2.6) {Nested Interval\\Property (NIP)};
        \node[thm] (CCC) at (0,-5.2) {Cauchy Convergence\\Criterion (CCC)};

        \draw[eqv] (CA) -- (MCT);
        \draw[eqv] (MCT) -- (BW);
        \draw[eqv] (BW) -- (HB);
        \draw[imp] ([xshift=-0.5cm]CA.south) -- node[lbl, left=2pt] {always} ([xshift=-0.5cm]NIP.north);
        \draw[imp] ([xshift=-0.5cm]NIP.south) -- node[lbl, left=2pt] {always} ([xshift=-0.5cm]CCC.north);
        \draw[imp, dashed] ([xshift=0.5cm]NIP.north) -- node[lbl, right=2pt] {+ Archimedean} ([xshift=0.5cm]CA.south);
        \draw[imp, dashed] (CCC.west) -- ++(-0.7,0) |- node[lbl, pos=0.25, rotate=90] {+ Archimedean} (CA.west);

        \node[font=\small, align=left, anchor=west] at (3.2,-3.9) {
            In any ordered field:\\
            $\bullet$ CA, MCT, BW, HB are equivalent, and each implies Archimedean;\\
            $\bullet$ CA $\Rightarrow$ NIP $\Rightarrow$ CCC, and neither converse holds in general;\\
            $\bullet$ the dashed arrows need the Archimedean property.
        };
    \end{tikzpicture}
    \end{center}
\end{remark}

